\documentclass[11pt]{article}

\usepackage[a4paper,margin=2.7cm]{geometry}
\usepackage{amsmath,amssymb,amsthm,mathtools}
\usepackage{enumitem}
\usepackage{hyperref}

\newtheorem{theorem}{Theorem}[section]
\newtheorem{proposition}[theorem]{Proposition}
\newtheorem{lemma}[theorem]{Lemma}
\newtheorem{remark}[theorem]{Remark}

\newcommand{\C}{\mathbb C}
\newcommand{\R}{\mathbb R}
\newcommand{\Z}{\mathbb Z}
\newcommand{\1}{\mathbf 1}
\newcommand{\dd}{\,d}
\newcommand{\pbar}{\partial_{\bar z}}
\newcommand{\pz}{\partial_z}
\newcommand{\calC}{\mathcal C}
\newcommand{\calN}{\mathcal N}

\title{Two One-Boundary Proofs of the Section 4 Rigidity Theorem}
\author{}
\date{}

\begin{document}
\maketitle

\section{Setup}

This note gives two self-contained ways to close the simply connected case in
Section~4 of \texttt{main.tex}.  The point is to use the fact that a simply
connected bounded domain has only one boundary component.  That removes the
inner-boundary cancellation obstruction which appears in the annular notes.

Let \(U\subset\C\) be bounded and simply connected.  Suppose that the Hermite
function \(h_k\) is an eigenfunction of the Gaussian localization operator with
symbol \(\1_U\).  As in Section~4 of \texttt{main.tex}, the Bargmann-side
identity implies
\begin{equation}\label{eq:moment-generating}
  \int_U e^{\pi\bar z w}|z|^{2k}e^{-\pi |z|^2}\dd A(z)=c
  \qquad(w\in\C)
\end{equation}
for a real constant \(c\).  Equivalently,
\begin{equation}\label{eq:bulk-moments}
  \int_U z^n |z|^{2k}e^{-\pi |z|^2}\dd A(z)=0,
  \qquad n\ge1,
\end{equation}
after conjugating the coefficients of \eqref{eq:moment-generating}.

The same divisibility argument used in Section~4 gives a compactly supported
potential \(u\) such that
\begin{equation}\label{eq:pde}
  \Delta u=|z|^{2k}e^{-\pi |z|^2}\1_U-c\,\delta_0
  \qquad\text{in }\mathcal D'(\C),
\end{equation}
and \(u\equiv0\) in the unbounded component of \(\C\setminus\overline U\).
Since \(U\) is simply connected, this is the whole exterior.  The singularity
forces \(0\in U\): if \(0\notin U\), then \(u\) would vanish near \(0\), which is
incompatible with the point mass in \eqref{eq:pde}; if \(0\in\partial U\), the
explicit logarithmic term contradicts \(u=0\) along exterior points tending to
\(0\).

We use the standard free-boundary regularity consequence of
\eqref{eq:pde}: because the exterior phase is identically zero and the forcing
\(|z|^{2k}e^{-\pi |z|^2}\) is real analytic and strictly positive near
\(\partial U\), the boundary \(\Gamma:=\partial U\) is a real-analytic Jordan
curve.  This is the usual Kinderlehrer--Nirenberg one-phase regularity input.
Consequently \(\Gamma\) has a Schwarz function \(S\) in a collar of the boundary:
\[
  S(z)=\bar z,\qquad z\in\Gamma .
\]

Set
\begin{equation}\label{eq:Phi-def}
  \Phi_k(t):=\int_0^t s^k e^{-\pi s}\dd s .
\end{equation}
Then \(\Phi_k\) is entire and strictly increasing on \([0,\infty)\).

\section{First Proof: The One-Boundary Moment Closure}

This proof is the clean version of the asymptotic route.  In the one-boundary
case one does not need a delicate saddle expansion: the vanishing of all
moments says exactly that the exterior Cauchy transform has zero expansion at
infinity.  The jump formula then produces a holomorphic interior function with
real boundary values.

\begin{proof}[One-boundary Cauchy-transform proof]
For \(n\ge1\),
\[
  \pbar\!\left(z^{n-1}\Phi_k(|z|^2)\right)
  =
  z^n |z|^{2k}e^{-\pi |z|^2}.
\]
Using \eqref{eq:bulk-moments} and Cauchy--Green,
\begin{equation}\label{eq:boundary-moments}
  \int_\Gamma z^{n-1}\Phi_k(|z|^2)\dd z=0,
  \qquad n\ge1.
\end{equation}
Equivalently,
\begin{equation}\label{eq:all-boundary-moments}
  \int_\Gamma z^m\Phi_k(|z|^2)\dd z=0,
  \qquad m\ge0.
\end{equation}

On the analytic collar of \(\Gamma\), define
\[
  Q(z):=\Phi_k(zS(z)).
\]
Then \(Q|_\Gamma=\Phi_k(|z|^2)\).  Consider the Cauchy transform
\[
  \calC_Q(\zeta):=\frac{1}{2\pi i}\int_\Gamma
  \frac{Q(z)}{z-\zeta}\dd z ,
  \qquad \zeta\notin\Gamma,
\]
where \(\Gamma\) is positively oriented.  For large \(|\zeta|\) in the exterior
component,
\[
  \calC_Q(\zeta)
  =
  -\sum_{m=0}^{\infty}
  \frac{1}{\zeta^{m+1}}
  \frac{1}{2\pi i}\int_\Gamma z^m Q(z)\dd z
  =0
\]
by \eqref{eq:all-boundary-moments}.  Since the exterior component is connected,
the identity theorem gives
\[
  \calC_Q\equiv0
  \qquad\text{in }\C\setminus\overline U .
\]

By the Plemelj jump relation,
\[
  \calC_Q^{\mathrm{int}}-\calC_Q^{\mathrm{ext}}=Q
  \qquad\text{on }\Gamma .
\]
Hence \(Q|_\Gamma\) is the boundary trace of the holomorphic function
\(\calC_Q^{\mathrm{int}}\) in \(U\).  But \(Q|_\Gamma=\Phi_k(|z|^2)\) is real-valued.
Therefore the imaginary part of \(\calC_Q^{\mathrm{int}}\) is harmonic in \(U\),
continuous up to \(\Gamma\), and zero on \(\Gamma\).  The maximum principle gives
\[
  \Im \calC_Q^{\mathrm{int}}\equiv0.
\]
Thus \(\calC_Q^{\mathrm{int}}\) is both holomorphic and real-valued, so it is
constant.  Restricting back to \(\Gamma\),
\[
  \Phi_k(|z|^2)=\kappa
  \qquad(z\in\Gamma)
\]
for some real constant \(\kappa\).  Since \(\Phi_k\) is strictly increasing on
\([0,\infty)\), \(|z|\) is constant on \(\Gamma\).  Hence \(\Gamma\) is a circle
centered at the origin, and the bounded simply connected domain \(U\) is the disk
bounded by that circle.
\end{proof}

\begin{remark}
This is the precise sense in which the ``asymptotic'' strategy becomes harmless
when there is only one boundary component.  The large-\(\zeta\) expansion of
\(\calC_Q\) is an exact asymptotic expansion at infinity, and all of its
coefficients vanish by \eqref{eq:all-boundary-moments}.  In a multiply connected
domain one instead gets a sum of component integrals, and different components
can in principle cancel.  That is the obstruction encountered in the annulus
notes.
\end{remark}

\section{Second Proof: The DtN Boundary Equation}

The second route packages the same rigidity as an overdetermined Dirichlet
problem.  It also gives the requested boundary parametrization equation.

Let \(\rho_k(r):=r^{2k}e^{-\pi r^2}\).  We keep the sign convention of
Section~4, namely \(\Delta u=\rho_k\1_U-c\delta_0\); replacing \(u\) by
\(-u\) gives the equivalent convention \(-\Delta u=\rho_k\1_U-c\delta_0\).
For boundary data \(\varphi\), let \(u_\varphi\) solve
\[
  \Delta u_\varphi=\rho_k(|z|)-c\,\delta_0
  \qquad\text{in }U,
  \qquad
  u_\varphi|_{\Gamma}=\varphi .
\]
The associated Dirichlet-to-Neumann map is
\[
  \calN_U(\varphi):=\partial_\nu u_\varphi|_{\Gamma}.
\]
The localization/free-boundary condition gives precisely
\[
  \calN_U(0)=0.
\]
Indeed, \(u\equiv0\) in the exterior and \(u\in C^1(\C\setminus\{0\})\), so on
\(\Gamma\) both \(u=0\) and \(\partial_\nu u=0\).  Since the tangential derivative
of a zero Dirichlet trace is also zero, this is the full condition
\[
  \nabla u=0
  \qquad\text{on }\Gamma .
\]

\begin{proof}[DtN proof and boundary ODE]
Define, for \(z\neq0\),
\[
  \Psi_k(z):=\frac{\Phi_k(|z|^2)}{z}.
\]
Since \(\Phi_k'(t)=t^k e^{-\pi t}\),
\[
  \pbar\Psi_k(z)=|z|^{2k}e^{-\pi |z|^2}
  \qquad(z\neq0),
\]
and the identity holds distributionally in \(U\) because \(\Phi_k(0)=0\).  Set
\[
  F(z):=4\pz u(z)-\Psi_k(z)+\frac{c}{\pi z},
  \qquad z\in U\setminus\{0\}.
\]
Using \(\Delta=4\pz\pbar\), \(\pbar(1/z)=\pi\delta_0\), and \eqref{eq:pde}, we get
\[
  \pbar F
  =
  \Delta u-\rho_k(|z|)+c\delta_0
  =0
  \qquad\text{in }\mathcal D'(U).
\]
Thus \(F\) is holomorphic in \(U\); the singularity at \(0\) is removable by the
same cancellation.

Now define
\[
  W(z):=\pi zF(z).
\]
Then \(W\) is holomorphic in \(U\).  Since \(\nabla u=0\) on \(\Gamma\), we have
\[
  W|_\Gamma
  =
  c-\pi\Phi_k(|z|^2),
\]
which is real-valued.  Hence \(\Im W\) is harmonic in \(U\), continuous up to
\(\Gamma\), and zero on \(\Gamma\).  By the maximum principle \(W\) is constant:
\[
  W\equiv\kappa .
\]
Therefore every real-analytic parametrization \(\gamma:\R/2\pi\Z\to\Gamma\)
satisfies the boundary equation
\begin{equation}\label{eq:param-equation}
  c-\pi\Phi_k(|\gamma(t)|^2)=\kappa .
\end{equation}
Differentiating gives
\[
  0
  =
  \frac{\dd}{\dd t}\Phi_k(|\gamma(t)|^2)
  =
  |\gamma(t)|^{2k}e^{-\pi|\gamma(t)|^2}
  \frac{\dd}{\dd t}|\gamma(t)|^2 .
\]
Since \(0\notin\Gamma\), the prefactor is strictly positive.  Hence
\begin{equation}\label{eq:param-ode}
  \frac{\dd}{\dd t}|\gamma(t)|^2=0,
  \qquad\text{equivalently}\qquad
  \Re\bigl(\overline{\gamma(t)}\,\gamma'(t)\bigr)=0.
\end{equation}
Thus \(\gamma'(t)=i a(t)\gamma(t)\) for a real function \(a(t)\), and
\(|\gamma(t)|\) is constant.  The boundary is a circle centered at the origin,
and simple connectivity gives that \(U\) is the disk it bounds.
\end{proof}

\begin{remark}
The DtN formulation also identifies the exact place where simple connectivity is
used.  With one boundary component, real boundary values of the holomorphic
function \(W\) force \(W\) to be constant.  In a multiply connected problem, one
would need the same real trace information on every boundary component, together
with the correct exterior/free-boundary interpretation on those components.  This
is precisely what the failed disk-annulus argument did not validly supply.
\end{remark}

\end{document}
